% Study note (not thesis text): the log-log scale, SCORE2 and the plaque model, built up slowly.
% Written 2026-11-14 on request ("explain it in the pdf very simply and slow build up").
% Sources: SCORE2 Updated Supplementary Material.docx (project root), read 2026-11-14: Box 1 (Fine and
%   Gray models), Box 6 (recalibration by intercept and slope), Supplementary Methods Tables 2-4
%   (baseline survival 0.9605/0.9776, scale1/scale2, worked example). Slope 0.722 from
%   CGPS_data/score2_honesty.md (person-level model, not the pooled GLMM).
% Numbers in Sections 2 and 5 recomputed in Python 2026-11-14.
\documentclass[11pt]{article}
\usepackage[a4paper,margin=25mm]{geometry}
\usepackage{amsmath,amssymb}
\usepackage{booktabs}
\usepackage{pgfplots}
\pgfplotsset{compat=1.18}
\usepackage[hidelinks]{hyperref}
\usepackage{parskip}
\usepackage{xcolor}
\usepackage{tcolorbox}
\newcommand{\up}{\textcolor{blue!70!black}}
\newcommand{\down}{\textcolor{red!70!black}}

\title{The log-log scale, SCORE2 and the plaque model\\[2pt]\large A slow explanation}
\author{Study note, 14 November 2026}
\date{}
\usepackage{titling}
\setlength{\droptitle}{-2.5cm}

\begin{document}
\maketitle

\section{The problem: a risk is stuck between 0 and 1}

Both SCORE2 and the plaque model predict a \emph{risk}: the probability $p$ that something happens
within a time window (10 years for SCORE2, the time between the two scans for plaque).
A risk is always between 0 and 1.

A regression model, on the other hand, builds a \emph{linear predictor}
\[
  \eta = \beta_0 + \beta_1 x_1 + \beta_2 x_2 + \dots
\]
which is a sum of terms and can be any number, from $-\infty$ to $+\infty$.

So we need a translation between the two: a way to turn a risk into a number that can be any
value, and back again. The log-log transform is one such translation, and SCORE2 and the plaque
model both use it.

\section{Going up: from risk to the log-log scale, one step at a time}

Start with a risk $p$. Apply three small steps.

\textbf{Step 1: flip to survival.}
\[
  S = 1 - p .
\]
$S$ is the probability that the event has \emph{not} happened by the end of the window.
Still between 0 and 1.

\noindent
\begin{minipage}[t]{0.38\textwidth}
\textbf{Step 2: take minus the log. This gives the cumulative hazard.}
\[
  \Lambda = -\log S = -\log(1 - p).
\]
The hazard is the rate at which the event happens at each moment; $\Lambda$ is that rate summed
over the whole window. It is between 0 and $+\infty$. One side is open now.

\medskip
See the red box for what ``minus the log'' does and why it means ``adding up risk''.
\end{minipage}\hfill
\begin{minipage}[t]{0.60\textwidth}
\vspace{-1.2em}
\begin{tcolorbox}[colback=red!4,colframe=red!70!black,boxrule=0.8pt,arc=2pt,
  left=5pt,right=5pt,top=4pt,bottom=4pt,fontupper=\footnotesize,
  title={\small What does ``take minus the log'' do?},fonttitle=\bfseries]
\textbf{Mechanically.} $S$ is between 0 and 1, and the log of such a number is always
negative: $\log 0.95 = -0.051$, $\log 0.60 = -0.511$. The minus flips it to positive. So
$-\log S$ runs from 0 ($S=1$, nobody has had the event) to $+\infty$ ($S=0$, everybody has).

\smallskip
\textbf{Meaning: it adds up risk over time.} Cut the window into years. To stay event-free for
10 years you must get through year 1, \emph{and} year 2, and so on, so the yearly survival
probabilities multiply:
\[
  S = 0.99^{10} = 0.904 .
\]
The log turns that product into a sum:
\[
  -\log S = 10 \times (-\log 0.99) = 10 \times 0.01005 = 0.1005 .
\]
So $-\log S$ is the small yearly risks added up: the hazard piled up over the window, which is
why $\Lambda$ is called \emph{cumulative}.

\smallskip
\textbf{Reading the number.} For small risks $\Lambda \approx p$ (risk 0.05 gives 0.051). For
large risks $\Lambda$ grows faster than $p$ (0.40 gives 0.51; 0.90 gives 2.30), as each extra
bit of risk is harder to reach when few are left event-free. $\Lambda = 1$ means
$S = e^{-1} = 0.37$.

\smallskip
\textbf{Why bother.} On the survival scale things combine by \emph{multiplying}; after minus
the log they combine by \emph{adding}, year 1 plus year 2, or segment 1 plus segment 2. That
is what lets the nine segment hazards sum into one person hazard.
\end{tcolorbox}
\end{minipage}
\medskip

\textbf{Step 3: take the log again.}
\[
  \eta = \log \Lambda = \log\{-\log(1-p)\}.
\]
Now the value can be anything from $-\infty$ to $+\infty$. Both sides are open, so this is the
number the linear predictor is allowed to be.

That is all ``log-log'' means: \textbf{log, then log again}. (The full name, complementary
log-log or cloglog, comes from the $1-p$ in step 1, the complement.)

\begin{table}[h]
\centering
\begin{tabular}{rrrr}
\toprule
risk $p$ & survival $S = 1-p$ & cum.\ hazard $\Lambda = -\log S$ & log-log $\eta = \log\Lambda$ \\
\midrule
0.01 & 0.99 & 0.0101 & $-4.60$ \\
0.05 & 0.95 & 0.0513 & $-2.97$ \\
0.10 & 0.90 & 0.1054 & $-2.25$ \\
0.40 & 0.60 & 0.5108 & $-0.67$ \\
0.90 & 0.10 & 2.3026 & $0.83$ \\
\bottomrule
\end{tabular}
\caption{The same five risks on each of the three scales. Notice that for small risks,
$\Lambda \approx p$.}
\end{table}

\section{Going down: from the log-log scale back to risk}

Every step can be undone. The opposite of $\log$ is $e^{(\cdot)}$, so going back means
using $e$ twice, in the reverse order:
\[
  \eta
  \;\xrightarrow{\ \ e^{\eta}\ \ }\;
  \Lambda
  \;\xrightarrow{\ \ e^{-\Lambda}\ \ }\;
  S
  \;\xrightarrow{\ \ 1 - S\ \ }\;
  p .
\]
Written in one line:
\[
  p = 1 - \exp\{-\exp(\eta)\}.
\]

So the rule of direction is simple:
\begin{center}
\begin{tabular}{ll}
\up{Going up} (risk $\to$ model scale): & $\log$, then $\log$ \quad $\eta = \log\{-\log(1-p)\}$ \\[3pt]
\down{Going down} (model scale $\to$ risk): & $\exp$, then $\exp$ \quad $p = 1-\exp\{-\exp(\eta)\}$
\end{tabular}
\end{center}
The model is \emph{fitted and added up} on the log-log scale (up). The \emph{answer you report}
is a risk (down).

\begin{figure}[h]
\centering
\begin{tikzpicture}
\begin{axis}[width=0.75\textwidth,height=6cm,xlabel={log-log value $\eta$},ylabel={risk $p$},
  xmin=-6,xmax=2,ymin=0,ymax=1,grid=major,grid style={gray!20}]
\addplot[very thick,blue!70!black,domain=-6:2,samples=200]{1-exp(-exp(x))};
\addplot[only marks,mark=*,red!70!black] coordinates {(-2.97,0.05) (-0.672,0.40)};
\node[anchor=west] at (axis cs:-2.8,0.09) {\small $p=0.05$, $\eta=-2.97$};
\node[anchor=west] at (axis cs:-0.5,0.36) {\small $p=0.40$, $\eta=-0.67$};
\end{axis}
\end{tikzpicture}
\caption{Going down: $p = 1-\exp\{-\exp(\eta)\}$. Any value of $\eta$ gives a risk between 0
and 1.}
\end{figure}

\section{Why log twice? Each log turns a multiplication into an addition}

The two logs are not arbitrary. Each one turns something that \emph{multiplies} into something
that \emph{adds}, and regression is built on adding.

\textbf{The first log} (survival $\to$ cumulative hazard).
Surviving two things means multiplying their survival probabilities:
$S = S_1 \times S_2$. After the minus-log, $\Lambda = \Lambda_1 + \Lambda_2$.
Hazards \emph{add}. This is why the nine segment hazards of one person can simply be added into
one person hazard.

\textbf{The second log} (cumulative hazard $\to$ log-log).
A risk factor such as smoking \emph{multiplies} the hazard, for example by 2:
$\Lambda_{\text{smoker}} = 2 \times \Lambda_{\text{non-smoker}}$.
After the log, $\log\Lambda_{\text{smoker}} = \log 2 + \log\Lambda_{\text{non-smoker}}$.
The effect \emph{adds} a fixed amount, $\beta = \log 2 \approx 0.69$. Going down again,
$e^{\beta} = 2$ is the hazard ratio.

Example: a non-smoker with risk 0.05 has $\Lambda = 0.0513$. A smoker with hazard ratio 2 has
$\Lambda = 0.1026$ and risk $1 - e^{-0.1026} = 0.098$. On the log-log scale this is just
$-2.97 + 0.69 = -2.28$.

So on the log-log scale every risk factor is one number that is added, and the same number
holds for everyone.

\section{What SCORE2 does}

From the SCORE2 Supplementary Material (Box 1 and Supplementary Methods Tables 2--4):

\textbf{The model.} SCORE2 was fitted with Fine and Gray models, a kind of proportional-hazards
model that also accounts for people dying of other causes first. It produces a linear predictor
LP from age, smoking, blood pressure, total and HDL cholesterol (and their interactions with
age).

\textbf{The uncalibrated risk} (Table 4, step 2):
\[
  \text{risk} = 1 - S_0^{\exp(\text{LP})},
\]
where $S_0$ is the baseline 10-year survival (0.9605 for men, 0.9776 for women).

Go \up{up} with the log-log steps from Section 2:
\[
  \log\{-\log(1-\text{risk})\} = \underbrace{\log\{-\log S_0\}}_{\text{a constant}} + \text{LP}.
\]
So on the log-log scale, SCORE2 is a \emph{straight line}: a constant plus the linear predictor.
SCORE2 lives on exactly the scale described above.

\textbf{The recalibration} (Table 4, step 3). To fit each European region, SCORE2 goes up to
the log-log scale, stretches and shifts the value with two numbers, and comes down again:
\[
  \text{calibrated risk} = 1 - \exp\Bigl[-\exp\bigl(\text{scale}_1 + \text{scale}_2 \times
  \log\{-\log(1 - \text{uncalibrated risk})\}\bigr)\Bigr].
\]
Read it from the inside out: \up{log, log} (up), multiply by $\text{scale}_2$ and add
$\text{scale}_1$ (a straight line), \down{exp, exp} (down).

\textbf{Worked example from Table 4} (man, 50 years, smoker, systolic BP 140, total cholesterol
6.3, HDL 1.4, low-risk region, as Denmark):
\begin{center}
\begin{tabular}{lll}
\toprule
step & calculation & value \\
\midrule
linear predictor & sum of the terms in Table 4 & $0.3221$ \\
uncalibrated risk & $1 - 0.9605^{\exp(0.3221)}$ & $0.0541$ \\
\up{up}: log-log & $\log\{-\log(1-0.0541)\}$ & $-2.889$ \\
straight line & $-0.5699 + 0.7476 \times (-2.889)$ & $-2.730$ \\
\down{down}: $\exp$ & $\exp(-2.730)$ & $0.0652$ \\
\down{down}: risk & $1 - \exp(-0.0652)$ & $0.0631$ \\
\bottomrule
\end{tabular}
\end{center}
The calibrated 10-year risk is 6.3\%, as in the supplement.

\section{What the plaque model does}

The plaque model works on the same scale. For segment $j$ of person $i$,
\[
  \log\{-\log(1 - p_{ij})\} = \eta_{ij} + b_i ,
\]
where $p_{ij}$ is the risk that the segment has plaque at follow-up, $\eta_{ij}$ holds the
covariates and $b_i$ is the person's own random level. Going \down{down} gives the segment
risk, and adding the nine segment hazards (Section 4) gives the person risk.

To use SCORE2 as the base, its risk is taken \up{up} to the log-log scale and put in as one
covariate:
\[
  \log\{-\log(1 - p_{ij})\} = \alpha_j + \gamma \times \log\{-\log(1 - \text{SCORE2})\}
  + \text{geometry terms} + b_i .
\]
Compare this with SCORE2's own recalibration in Section 5:
$\alpha$ plays the role of $\text{scale}_1$ and $\gamma$ the role of $\text{scale}_2$.
\textbf{The plaque model recalibrates SCORE2 to plaque in exactly the way SCORE2 recalibrates
itself to each region.}

\section{Why it is good to do it on the same scale}

\textbf{The slope has a clear meaning.} Both sides are log cumulative hazards, so $\gamma$
compares like with like.
If $\gamma = 1$, a person whose SCORE2 hazard is twice as high also has a twice as high plaque
hazard: SCORE2's differences carry over unchanged.
If $\gamma < 1$, the differences shrink: twice the SCORE2 hazard means $2^{\gamma}$ times the
plaque hazard. Your person-level estimate (\texttt{score2\_honesty.md}) is $\gamma = 0.722$
(95\% CI 0.569 to 0.874), so twice the SCORE2 hazard gives about $2^{0.722} \approx 1.65$ times
the plaque hazard. SCORE2 spreads people out a bit more than their plaque risk differs.

\textbf{A straight line is enough.} SCORE2 is a straight line on the log-log scale, so a
straight line (two numbers, $\alpha$ and $\gamma$) is the natural way to move it to a new
outcome. Putting SCORE2 in as a raw percentage, or on the logit scale, would bend that
relation, and the slope would no longer mean ``how much of SCORE2's hazard difference carries
over''.

\textbf{It follows the published method.} The intercept-and-slope recalibration is the one the
SCORE2 authors use themselves (Box 6), so it can be cited directly.

\textbf{Geometry adds on the same scale.} The geometry coefficients are added on the same
log-log scale as SCORE2, so $e^{\beta}$ is a hazard ratio on top of SCORE2: how many times
higher the plaque hazard is for one unit of geometry, at the same SCORE2 value.

\section{What to be careful about}

\begin{itemize}
\item \textbf{A different kind of hazard.} SCORE2 is a Fine and Gray model, so its hazard is a
  \emph{subdistribution} hazard, which accounts for death from other causes. The supplement
  itself says its hazard ratios ``are relevant for risk estimation only''. The plaque model has
  an ordinary hazard. The scale is the same in form, but the two hazards are not the same thing.
\item \textbf{A different outcome.} SCORE2 predicts cardiovascular events, not plaque. A
  $\gamma$ below 1 is expected for that reason alone.
\item \textbf{A different model.} The 0.722 comes from the person-level model in
  \texttt{score2\_honesty.md}. The decision on the base must be made in the pooled segment
  model.
\item \textbf{Which SCORE2 column.} The CGPS \texttt{score2} value is most likely the
  region-calibrated risk. That does not change the argument: the calibration is itself a
  straight line on the log-log scale, so the value is still a straight line in LP.
\end{itemize}

\end{document}
